\chapter{Shingles (Herpes Zoster)}

\section*{Definition}
Shingles is a painful, blistering rash caused by reactivation of the varicella-zoster virus — the same virus that causes chickenpox. It usually appears as a stripe of blisters on one side of the body, most often on the torso or face.

\section*{What Happens in Your Body}
After you recover from chickenpox, the virus stays in your body in a dormant (inactive) state inside nerve roots near your spinal cord. Years or decades later, the virus can reactivate, often when your immune system is weakened. The virus travels along the nerve to your skin, causing inflammation, pain, and a blistering rash in the area served by that nerve. Because each nerve supplies a specific strip of skin, the rash typically stays on one side of your body in a band-like pattern.

\section*{Causes}
\begin{itemize}
  \item Reactivation of the varicella-zoster virus (chickenpox virus) — anyone who has had chickenpox is at risk
  \item Aging — risk increases sharply after age 50
  \item Weakened immune system: Chemotherapy, organ transplant, HIV, long-term steroid use, or immune-suppressing medications
  \item Physical or emotional stress
  \item Serious illness or injury
\end{itemize}

\section*{When to See a Doctor}
See your doctor promptly if:
\begin{itemize}
  \item You develop a painful, blistering rash on one side of your body — early treatment within 72 hours reduces severity and risk of complications
  \item The rash is on your face, near your eye, or on your nose — shingles near the eye can threaten vision
  \item You have severe pain that is not controlled by over-the-counter pain relievers
  \item You are over 60 and develop shingles (higher risk of complications)
  \item You have a weakened immune system
  \item The rash spreads widely or looks like it might be infected (warm, red, oozing pus)
  \item Pain persists after the rash heals (postherpetic neuralgia)
\end{itemize}

\section*{Related Symptoms}
\begin{itemize}
  \item Burning, tingling, or shooting pain before the rash appears (prodromal pain)
  \item Red, fluid-filled blisters that crust over within 7-10 days
  \item Itching or sensitivity in the affected area
  \item Fever and chills
  \item Headache
  \item Fatigue
  \item Sensitivity to light
\end{itemize}
\textbf{Important:} Shingles is contagious to people who have never had chickenpox or the chickenpox vaccine — they can catch chickenpox from contact with your open blisters. Keep the rash covered until all blisters have crusted over.

\section*{Self-Care / Home Management}
\begin{itemize}
  \item Start antiviral medication within 72 hours of the rash appearing — this reduces pain, shortens the illness, and lowers the risk of long-term nerve pain.
  \item Keep the rash clean and dry to prevent bacterial infection.
  \item Wear loose, soft clothing over the affected area to reduce irritation.
  \item Apply cool, wet compresses to blisters for 20 minutes several times a day to soothe pain and itching.
  \item Take acetaminophen (Tylenol) or ibuprofen (Advil, Motrin) for pain relief.
  \item Avoid contact with pregnant women, newborns, and people with weakened immune systems until blisters have fully crusted over.
\end{itemize}

\section*{How Your Doctor Will Evaluate You}
\begin{itemize}
  \item Your doctor can usually diagnose shingles by looking at the rash and asking about your symptoms — the one-sided, band-like pattern is characteristic.
  \item They may swab a blister to test for the virus if the diagnosis is uncertain.
  \item If shingles involves your eye, you will be referred to an eye specialist (ophthalmologist).
  \item Your doctor will ask about your medical history, immune status, and whether you have had chickenpox or the shingles vaccine.
\end{itemize}

\section*{Treatment Options}
\begin{itemize}
  \item Antiviral medications (acyclovir, valacyclovir, famciclovir) — most effective when started within 72 hours of rash onset; they reduce pain, speed healing, and lower the risk of complications.
  \item Pain medications — over-the-counter options for mild pain; prescription nerve pain medications (gabapentin, pregabalin) or lidocaine patches for more severe pain.
  \item Corticosteroids may be used in some cases to reduce inflammation, especially when the rash is severe or involves the eye.
  \item The shingles vaccine (Shingrix) is recommended for adults 50 and older and for younger adults with weakened immune systems — it reduces the risk of getting shingles by more than 90\% and also lowers the chance of long-term nerve pain.
\end{itemize}

\section*{References}
\begin{thebibliography}{99}
\bibitem{ref1} Cohen JI. "Clinical practice: Herpes zoster." \textit{New England Journal of Medicine}, 2013;369(3):255--263.
\bibitem{ref2} Oxman MN, Levin MJ, Johnson GR, et al. "A vaccine to prevent herpes zoster and postherpetic neuralgia in older adults." \textit{New England Journal of Medicine}, 2005;352(22):2271--2284.
\bibitem{ref3} Johnson RW, Rice ASC. "Postherpetic neuralgia." \textit{New England Journal of Medicine}, 2014;371(16):1526--1533.
\bibitem{ref4} Arvin AM. "Varicella-zoster virus." \textit{Clinical Microbiology Reviews}, 1996;9(3):361--381.
\bibitem{ref5} Dworkin RH, Johnson RW, Breuer J, et al. "Recommendations for the management of herpes zoster." \textit{Clinical Infectious Diseases}, 2007;44(Suppl 1):S1--S26.
\bibitem{ref6} Cunningham AL, Lal H, Kovac M, et al. "Efficacy of the adjuvanted recombinant zoster vaccine in adults aged 50 years or older." \textit{New England Journal of Medicine}, 2016;375(11):1019--1032.
\end{thebibliography}
